% COMPILATION:
% pdflatex bvmt_short_ieee_article.tex (twice for cross-references)
% Target: <= 10 pages, IEEE conference format.
% Companion to the long version bvmt_multiagent_ieee_article.tex; same
% empirical content compressed to the IEEE conference page budget.

\documentclass[conference]{IEEEtran}
\IEEEoverridecommandlockouts

\usepackage{cite}
\usepackage{amsmath,amssymb,amsfonts}
\usepackage{graphicx}
\usepackage{textcomp}
\usepackage{xcolor}
\usepackage{booktabs}
\usepackage{multirow}
\usepackage{array}
\usepackage{url}
\usepackage{hyperref}
\usepackage{balance}

\begin{document}

\title{Reproducibility-First Multi-Component Forecasting for Illiquid Equity Markets: Baselines, Ablations, and a Documented Discrepancy on the Tunisian Stock Exchange}

\author{
\IEEEauthorblockN{Aymen Ben Brik}
\IEEEauthorblockA{\textit{Department of Information Technology and Mathematics}\\
\textit{Esprit School of Business}\\
Tunis, Tunisia\\
aymen.benbrik@esprit.tn}
\and
\IEEEauthorblockN{Amen Allah Fraj}
\IEEEauthorblockA{\textit{Department of Information Technology and Mathematics}\\
\textit{Esprit School of Business}\\
Tunis, Tunisia\\
amenallah.fraj@esprit.tn}
}

\maketitle

\begin{abstract}
We present a reproducibility study of multi-component deep forecasting on a thin emerging market (BVMT, 68 stocks, 2016--2025, 116{,}626 trading records). The pipeline combines a Temporal Fusion Transformer (TFT), XGBoost fundamentals, French NLP sentiment, a Graph Attention Network, and SHAP-based explanations. On a harmonized 2025 test slice ($N=5{,}582$, UP ratio $49.84\%$), the Vertex-reproducible TFT v3 reference reaches $51.79\%$ directional accuracy (95\% bootstrap CI $[50.43, 53.08]$, macro-F1 $0.516$), statistically indistinguishable from ten classical and deep baselines (logistic regression $52.79\%$, LSTM $52.69\%$, standard Transformer $52.45\%$; pairwise McNemar $p \in \{0.35, 0.41, 0.52\}$). Two internal TFT ablations move accuracy by $\leq 0.4$~pp ($p > 0.65$); the per-stock target normalizer and the Variable Selection Network are not measurably helpful on this dataset. An earlier internal estimate of $77.4\%$ is documented as non-reproducible from released artefacts and retained transparently as a development-time number, not as a result. A multi-agent ablation evaluated on offline-realizable configurations (TFT alone, TFT $+$ Graph, $\pm$ confidence gating) shows the graph signal cannot flip the binary technical decision under the deployed orchestrator weights; six remaining configurations are DB-dependent and listed as future work. The contributions are (i) a reproducible Vertex-AI training pipeline with two TFT ablations, (ii) ten baselines plus three TFT variants harmonized on a single test slice with bootstrap CIs and pairwise McNemar tests, and (iii) a documented gap between an earlier headline and the reproducible numbers, offered as a reproducibility lesson for thin-market AI research.
\end{abstract}

\begin{IEEEkeywords}
Temporal Fusion Transformers, emerging markets, explainable AI, financial forecasting, reproducibility, thin liquidity, EU AI Act-aligned transparency
\end{IEEEkeywords}

\section{Introduction}

Financial AI is overwhelmingly designed on liquid developed markets and breaks on thin emerging markets because of price-continuity gaps, heterogeneous per-stock scales, no derivatives, multilingual news (French + Arabic on BVMT), and no public benchmarks. The Tunisian Stock Exchange (BVMT, 68 listed securities, daily volume $0.5$--$2$ MDT) is a representative testbed. To our knowledge --- after IEEE Xplore, ACM Digital Library, arXiv, SSRN, and Google Scholar restricted to \{BVMT, Tunis Stock Exchange, North African equity forecasting\} (January 2026) --- no peer-reviewed deep-learning system has been published on BVMT-listed equities.

This paper centers reproducibility. We release a containerized Vertex-AI training pipeline that re-trains a TFT v3 reference (per-stock \texttt{GroupNormalizer}, Variable Selection Network, pinball quantile loss on $Q_{0.1}/Q_{0.5}/Q_{0.9}$, hidden $64$, encoder $30$, horizon $7$) plus two internal training ablations and one post-hoc evaluation-time ablation, bit-for-bit from the same seed. We harmonize evaluation across all $13$ models tested in this work (ten baselines plus three TFT variants) on a single 2025 test slice with bootstrap CIs and pairwise McNemar tests, document a non-reproducible earlier estimate of $77.4\%$ directional accuracy that does not survive a Vertex retrain, and report a partial multi-agent ablation whose offline-realizable subset shows the graph signal carries no detectable directional skill above TFT v3 alone.

The contribution is methodological rather than headline-accuracy. Our reproducible TFT v3 sits in the $50$--$53\%$ band shared with classical baselines and below the $55$--$65\%$ developed-market TFT range~\cite{IEEEhowto:kopka}, consistent with the thin-market priors of Bekaert and Harvey~\cite{ref:bekaert_harvey,ref:harvey_multiple,ref:harvey_presidential}.

\section{Related Work}\label{sec:related}

\textbf{Temporal models.} LSTMs~\cite{ref:lstm} suffer representation bottlenecks on multivariate heterogeneous inputs. Attention mechanisms~\cite{ref:vaswani} and the Temporal Fusion Transformer~\cite{IEEEhowto:kopka} add a Variable Selection Network that learns per-feature importance before temporal processing. Recent extensions optimize long-horizon efficiency~\cite{ref:informer,ref:autoformer,ref:fedformer,ref:patchtst,ref:itransformer} or use patch/frequency decompositions~\cite{ref:timesnet,ref:timemixer}; simple linear and hierarchical baselines~\cite{ref:dlinear,ref:nhits} match or exceed transformer accuracy on standard benchmarks, motivating careful baselining. We retain TFT for its Variable Selection Network's per-variable interpretability, which an EU AI Act-aligned audit trail can consume directly. Most published TFT applications target developed markets with directional accuracy $55$--$65\%$ on daily forecasts; on BVMT 2025 our Vertex-reproducible TFT v3 reaches $51.79\%$, below this range and consistent with the literature prior on emerging-market unpredictability~\cite{ref:bekaert_harvey,ref:harvey_multiple}.

\textbf{Multi-agent systems.} Classical autonomous Multi-Agent Systems~\cite{ref:multi_agent} feature heterogeneous learning agents trading on a shared environment; recent LLM-based frameworks~\cite{ref:yang_fingpt,ref:lopez_lira_chatgpt} use large language models as analyst/trader/risk-officer personas. Our system belongs to neither tradition: the seven components are specialized deterministic Python modules invoked by a central orchestrator, with typed structured outputs and no peer-to-peer communication. We frame it as a \emph{multi-component decision pipeline with explicit trust boundaries}, retaining the multi-agent vocabulary because of (i) the typed \texttt{AgentSignal} interface, (ii) centralized failure isolation~\cite{ref:ensemble}, and (iii) per-component auditability.

\textbf{Explainable AI and EU AI Act.} Articles 13 and 14 of the EU AI Act~\cite{ref:eu_ai_act} establish transparency and human-oversight requirements for high-risk AI in finance. Explainability~\cite{ref:explainability} requires faithfulness, contrastiveness, and selectivity. SHAP~\cite{ref:shap} provides a Shapley-value attribution, with known limitations under adversarial perturbation~\cite{ref:slack_fooling,ref:alvarez_melis,ref:kumar_shap} and unresolved gaps relative to causal explanation~\cite{ref:veale_aiact,ref:mokander_audit}. We integrate SHAP for XGBoost and attention weights for TFT; the resulting decision traces are designed to \emph{support} EU AI Act transparency, not by themselves to establish conformity.

\textbf{GNNs, NLP, and reproducibility.} Stock-network designs combine GAT-style attention~\cite{ref:gat,ref:fingat,ref:mansf}; we use a 60-day rolling correlation graph with a $0.40$ threshold. Domain-specific French NLP outperforms multilingual baselines on financial sentiment~\cite{ref:financial_nlp}; AraBERT~\cite{ref:arabert} is the standard Arabic-language baseline. We integrate the reproducibility-by-checklist recommendations of~\cite{ref:pineau_repro,ref:bouthillier_repro,ref:lipton_troubling}, document the backtest-overfitting risk known in quantitative finance~\cite{ref:bailey_pseudo,ref:harvey_presidential,ref:deflated_sharpe}, and use the strictly sequential walk-forward protocol of~\cite{ref:bergmeir_cv_ts}.

\section{System and Methodology}\label{sec:methodology}

\subsection{Multi-Component Pipeline}

The pipeline implements seven specialized modules orchestrated by a central coordinator (Fig.~\ref{fig:architecture_compact}). Each agent consumes a standardized input (the shared \texttt{AgentState}) and produces a typed \texttt{AgentSignal} dataclass containing a numerical signal, a calibrated confidence in $[0,1]$, an explanation dictionary, and an error flag. No agent has direct write access to another's state.

\begin{figure}[htbp]
\centering
\fbox{\parbox{0.92\columnwidth}{\centering\small
Stock input $\to$ DataAgent $\to$ TechnicalAgent (TFT) \\
\hspace*{1.5em}$\to$ FundamentalAgent (XGBoost) $\to$ SentimentAgent (NLP) \\
\hspace*{1.5em}$\to$ GraphAgent (GAT) $\to$ XAIAgent (SHAP)\\
$\Downarrow$\\
OrchestratorAgent (adaptive weights, confidence gating, failure isolation)\\
$\Downarrow$\\
Final signal + explanation}}
\caption{Seven-component pipeline architecture. Agents communicate through typed \texttt{AgentSignal} objects; the orchestrator adjusts weights, applies confidence gating, and redistributes weights on agent failure.}
\label{fig:architecture_compact}
\end{figure}

The orchestrator starts from a fixed base weight vector $w^{(0)} = (w_\text{tech}, w_\text{fund}, w_\text{sent}, w_\text{graph}) = (0.34, 0.26, 0.20, 0.20)$ and mutates it at request time according to data-quality conditions: empty news $\to$ sentiment weight reduced $80\%$ and redistributed; missing recent financials $\to$ fundamental weight reduced $50\%$; price history $<252$ rows $\to$ technical weight reduced $40\%$. Final weights are renormalized to sum to $1$. A confidence gate applies $w_\text{gated} = w_\text{orig} \cdot (\text{confidence}/\text{threshold})$ when an agent's confidence falls below an agent-specific threshold (technical $0.60$, fundamental $0.55$, sentiment $0.50$, graph $0.45$). Failure of agent $j$ triggers redistribution $\hat{w}_i = w_i / (1 - w_j)$ among the remaining agents.

\subsection{Dataset and Walk-Forward Protocol}

The dataset (Table~\ref{tab:dataset_compact}) combines daily prices, computed technical indicators, news articles, financial statements, and ratios from PostgreSQL tables.

\begin{table}[htbp]
\centering
\caption{Dataset (post-cleaning): TFT features sourced from cleaned daily-prices joins; train 2016--2023, val 2024, test 2025; \texttt{tft\_features.csv} is the public release.}
\label{tab:dataset_compact}
\footnotesize
\setlength{\tabcolsep}{4pt}
\begin{tabular}{lccc}
\toprule
\textbf{Split} & \textbf{Rows} & \textbf{Period} & \textbf{Tickers} \\
\midrule
Train               & 84{,}648  & 2016--2023 & 68 \\
Val                 & 10{,}228  & 2024       & 67 \\
Test (raw)          & 10{,}314  & 2025-01-02--2025-12-22 & 43 \\
Test (harmonized)   & 5{,}582   & 2025-01-07--2025-12-10 & 43 \\
\midrule
News articles       & 7{,}937   & 2016--2026 & 68 \\
Financial ratios    & 261       & 2016--2024 & 68 \\
\bottomrule
\end{tabular}
\end{table}

Walk-forward validation: $\mathcal{D}_\text{train} = \{\text{date} < 2024\text{-}01\text{-}01\}$, $\mathcal{D}_\text{val} = \{2024\}$, $\mathcal{D}_\text{test} = \{2025\}$. Random train/test splits are invalid for time series~\cite{ref:bergmeir_cv_ts}.

\subsection{TFT v3 Configuration}

The TechnicalAgent uses 15 engineered features (5 price, 4 momentum, 3 volatility, 3 derived relative features) on a 30-day encoder window, predicting the 7-day fractional return via pinball quantile loss on $Q_{0.1}/Q_{0.5}/Q_{0.9}$ (clipped to $[-15\%, +15\%]$). The Variable Selection Network learns per-variable softmax-gated importance weights $\boldsymbol{\xi}_t = \text{softmax}(W_\text{vsn}\,\mathbf{x}_t + \mathbf{b}_\text{vsn})$ before temporal processing. Two normalization mechanisms are causal by construction:
\begin{itemize}
\item \emph{Target}: \texttt{GroupNormalizer} fits $(\mu_i, \sigma_i)$ per stock $i$ on the training split only and applies them frozen to val/test, preventing cross-stock leakage and isolating each stock's return distribution.
\item \emph{Input features}: \texttt{EncoderNormalizer} refits a local mean and standard deviation \emph{inside} each 30-day encoder window at iteration time. Window-local statistics are causal because the window only contains information available at the prediction time.
\end{itemize}
Hyperparameters are summarized in Table~\ref{tab:tft_config}; the released Vertex container converges in $20$--$37$ epochs ($71$--$161$~min on a T4, $\sim\!\$8$ total).

\begin{table}[htbp]
\centering
\caption{TFT v3 hyperparameter configuration (\texttt{training/train\_tft\_v3.py}). Same settings are reused for the two internal ablations of Section~\ref{sec:tft_ablations}.}
\label{tab:tft_config}
\scriptsize
\setlength{\tabcolsep}{4pt}
\begin{tabular}{lc}
\toprule
\textbf{Hyperparameter} & \textbf{Value} \\
\midrule
hidden size               & $64$               \\
attention heads           & $4$                \\
dropout                   & $0.3$              \\
encoder length            & $30$ trading days  \\
prediction horizon        & $7$ trading days   \\
quantiles                 & $\{0.10, 0.50, 0.90\}$ \\
learning rate             & $3 \times 10^{-4}$ \\
warmup epochs             & $3$ (linear $10\% \to 100\%$) \\
optimizer                 & AdamW, $\ell_2 = 10^{-4}$ \\
gradient clip             & $0.5$              \\
max epochs / patience     & $50$ / $8$         \\
batch size / precision    & $128$ / bf16 mixed \\
target normalizer         & \texttt{GroupNormalizer(groups=[ticker\_id])} \\
input normalizer          & \texttt{EncoderNormalizer} (per-window)    \\
\bottomrule
\end{tabular}
\end{table}

The base agent weights of the orchestrator are $w^{(0)} = (0.34, 0.26, 0.20, 0.20)$ for (technical, fundamental, sentiment, graph). Three data-quality conditions rescale the vector at request time (Table~\ref{tab:weight_adjustments}); after all triggered rules, weights are renormalised to $\sum_a w_a = 1$. These heuristics were chosen during development against held-out 2024 data and have not been formally tuned by cross-validation; the offline-realisable subset of the multi-agent ablation evaluates whether they carry measurable directional skill.

\begin{table}[htbp]
\centering
\caption{Heuristic weight adjustments applied by the orchestrator's \texttt{\_adjust\_weights} based on per-ticker data quality. Each rule decreases the donor weight by a fixed fraction $\alpha$ and redistributes the recovered mass over the remaining agents in the proportions shown.}
\label{tab:weight_adjustments}
\scriptsize
\setlength{\tabcolsep}{3pt}
\begin{tabular}{p{2.3cm}cccc}
\toprule
\textbf{Trigger} & \textbf{Donor ($\alpha$)} & \textbf{Tech} & \textbf{Fund} & \textbf{Sent} \\
\midrule
\texttt{news\_count == 0}              & sent ($0.80$) & $+0.50$ & $+0.30$ & ---     \\
$1 \le$ news $< 5$                     & sent ($0.55$) & $+0.50$ & $+0.30$ & ---     \\
$\neg$ recent\_financials              & fund ($0.50$) & $+0.70$ & ---     & ---     \\
price\_rows $< 252$                    & tech ($0.40$) & ---     & $+0.35$ & $+0.35$ \\
\bottomrule
\end{tabular}
\end{table}

The quantile-to-direction rule is the median sign: $\hat{d}(x) = \text{UP}$ iff $\hat{Q}_{0.5}(x) \geq 0$. Confidence is read from the central interval width $c(x) = \text{clip}(1 - (\hat{Q}_{0.9} - \hat{Q}_{0.1})/\Delta_\text{ref}, 0.05, 0.99)$ with $\Delta_\text{ref} = 0.10$, with a multiplicative penalty of $0.6$ when the quantiles cross.

\textbf{Alternative decision rules.} The median-sign rule discards information. Three natural alternatives are: (\textit{R2}) strict bullish/bearish gating --- $\hat{d}\!=\!\text{UP}$ iff $\hat{Q}_{0.1}\!>\!0$, $\text{DOWN}$ iff $\hat{Q}_{0.9}\!<\!0$, otherwise abstain (higher precision, lower coverage); (\textit{R3}) HOLD band --- $\hat{d}\!=\!\text{HOLD}$ iff $|\hat{Q}_{0.5}|\!<\!\tau$ for a validation-tuned $\tau$; (\textit{R4}) probability via interpolated CDF --- treat $\{(0.1, \hat{Q}_{0.1}), (0.5, \hat{Q}_{0.5}), (0.9, \hat{Q}_{0.9})\}$ as a piecewise-linear inverse CDF, recover $\hat{p}(\text{return}\!>\!0)$, threshold at $0.5$. The Vertex-reproducible per-row CSV ships all three quantiles so future readers can re-evaluate accuracy under R2--R4 without re-training; this paper reports only R1 throughout.

\textbf{TFT version progression.} Table~\ref{tab:tft_progression} traces the validation loss from v1 (BCE on binary direction) through the non-reproducible v3 development run to the Vertex-reproducible v3 reference. The three rows live on different loss scales; the development pinball value $0.5161$ is measured on raw clipped fractional return, while the Vertex pinball value $0.0159$ is measured after the \texttt{GroupNormalizer}'s z-scoring, so the two are not directly comparable in absolute terms. The Vertex retrain converges early (epoch 4 best) under the same hyperparameters, suggesting the longer convergence trajectory of the development run was an artefact of the pre-Vertex pipeline whose exact cause we cannot recover.

\begin{table}[htbp]
\centering
\caption{TFT validation-loss progression. Three loss scales are reported: BCE (v1, binary direction target), pinball on raw clipped fractional return (v3 development, non-reproducible), and pinball on z-scored target (v3 Vertex, reproducible reference).}
\label{tab:tft_progression}
\scriptsize
\setlength{\tabcolsep}{4pt}
\begin{tabular}{lccc}
\toprule
\textbf{Version} & \textbf{Loss type} & \textbf{Val loss} & \textbf{Best epoch} \\
\midrule
Random baseline               & BCE                   & $0.6931$ & ---        \\
TFT v1 (consolidated)         & BCE                   & $0.6051$ & $7$        \\
TFT v3 dev (non-reproducible) & Pinball (raw return)  & $0.5161$ & $16$       \\
\textbf{TFT v3 Vertex (ref.)} & Pinball (z-scored)    & $\mathbf{0.0159}$ & $\mathbf{4}$ \\
\bottomrule
\end{tabular}
\end{table}

\section{Experimental Evaluation}\label{sec:experiments}

\subsection{Baselines and Harmonization}

We evaluate ten competing baselines (Table~\ref{tab:baselines}) on a shared walk-forward protocol: five trivial / trend-following rules (always-UP, always-DOWN, momentum $5/20$, MA$_{20}$ crossover), three tabular-ML models (logistic regression L2, random forest, XGBoost) trained on a flattened 60-day $\times$ 15-feature input, and two sequential-deep baselines (a 2-layer LSTM at $54$\,k parameters, a standard 2-block Transformer at $68$\,k parameters). All ten plus three TFT variants are evaluated on a single harmonized 2025 test slice ($N=5{,}582$, UP ratio $49.84\%$) obtained by intersecting the baseline grid with the de-duplicated Vertex grid (one valid prediction per (ticker, date), Section~\ref{sec:directional_acc}). All CIs are 95\% percentile-bootstrap (2{,}000 resamples).

\textbf{CNN-LSTM preliminary observation.} A preliminary CNN-LSTM ablation across four independent training runs with varied batch sizes and learning rates returns test accuracies in the band $51.7$--$53.8\%$, in the same $50$--$54\%$ regime as the rest of the panel. We are explicit that four runs do not power a formal ceiling proof: with a few-thousand-sample test set, the per-run binary-accuracy standard deviation is typically $0.5$--$1.0$ pp, comparable to the observed $2.1$ pp range. A seed-sweep harness ($N \geq 20$ seeds with mean $\pm 95\%$ bootstrap CI and Wilcoxon test) is shipped as \texttt{training/cnn\_lstm\_seed\_sweep.py}; the four-run observation is reported as a consistent signature, not an empirical ceiling. The hypothesised mechanism --- that $\text{Conv1D}$ over heterogeneous financial features (price, RSI, volume) treats them as spatially equivalent and lacks the per-feature gating that the TFT's Variable Selection Network provides --- is offered as the most parsimonious explanation pending the formal sweep.

\begin{table*}[htbp]
\centering
\caption{Directional accuracy of ten classical/deep baselines plus three TFT variants on the harmonized 2025 test slice ($N=5{,}582$, UP $49.84\%$). Pairwise McNemar two-sided $p$-values vs.\ TFT v3.}
\label{tab:baselines}
\small
\setlength{\tabcolsep}{4pt}
\begin{tabular}{llcccc}
\toprule
\textbf{Family} & \textbf{Model} & \textbf{Acc} & \textbf{95\% CI} & \textbf{Macro-F1} & \textbf{McN.\ vs TFT v3} \\
\midrule
\multirow{4}{*}{Trivial / trend}
& Always-UP                                       & $49.84\%$ & $[48.58, 51.15]$ & $0.333$ & $p = 0.054$ \\
& Always-DOWN                                     & $50.16\%$ & $[48.89, 51.45]$ & $0.334$ & $p = 0.069$ \\
& Momentum 5/20                                   & $50.36\%$ & $[49.10, 51.65]$ & $0.502$ & $p = 0.127$ \\
& MA$_{20}$ crossover                             & $49.77\%$ & $[48.50, 51.09]$ & $0.496$ & $p = 0.031$ \\
\midrule
\multirow{3}{*}{Tabular ML}
& Logistic regression (L2)                        & $52.79\%$ & $[51.45, 54.14]$ & $\mathbf{0.524}$ & $p = 0.351$ \\
& Random forest ($n_\text{trees}{=}200$)         & $50.38\%$ & $[49.01, 51.67]$ & $0.450$ & $p = 0.154$ \\
& XGBoost (400 trees)                             & $52.15\%$ & $[50.86, 53.49]$ & $0.513$ & $p = 0.749$ \\
\midrule
\multirow{2}{*}{Deep baselines}
& Simple LSTM (54 k params)                       & $52.69\%$ & $[51.33, 53.94]$ & $0.517$ & $p = 0.414$ \\
& Standard Transformer (68 k params)              & $52.45\%$ & $[51.15, 53.69]$ & $0.489$ & $p = 0.517$ \\
\midrule
\multirow{3}{*}{TFT (this work)}
& \textbf{TFT v3 reference} (Vertex)              & $\mathbf{51.79\%}$ & $[50.43, 53.08]$ & $0.516$ & --- \\
& \quad $-$ VSN (uniform weights)                  & $51.56\%$ & $[50.23, 52.97]$ & $0.503$ & $p = 0.734$ \\
& \quad $-$ per-stock norm                         & $51.40\%$ & $[50.07, 52.69]$ & $0.471$ & $p = 0.672$ \\
\bottomrule
\end{tabular}
\end{table*}

\subsection{Directional Accuracy on Test Data}\label{sec:directional_acc}

\textbf{Reproducible TFT v3 on the harmonized slice.} The Vertex retrain of the reference architecture (\texttt{vertex/Dockerfile} + \texttt{vertex/cloudbuild.yaml}) produces \texttt{baseline\_test\_predictions.csv}. The raw output ($6{,}782$ rows) contains duplicate inference outputs for some (ticker, date) pairs because the multi-horizon decoder writes one row per horizon step; after de-duplicating to one valid prediction per (ticker, date) we obtain the canonical slice of $\mathbf{N=5{,}582}$ predictions over $43$ tickers (UP ratio $\mathbf{49.84\%}$). TFT v3 reaches $\mathbf{51.79\%}$ directional accuracy (95\% bootstrap CI $[50.43, 53.08]$, macro-F1 $0.516$). It is statistically indistinguishable from the strongest non-TFT baselines (logistic regression $52.79\%$, $p = 0.35$; simple LSTM $52.69\%$, $p = 0.41$; standard Transformer $52.45\%$, $p = 0.52$) and from always-UP on the same slice ($49.84\%$, $p = 0.054$).

\smallskip
\noindent\fbox{\parbox{0.93\columnwidth}{\small \textbf{Reproducibility transparency.} An earlier internal training run was reported as reaching $\mathbf{77.4\%}$ directional accuracy on the broader 2025 test set ($N=10{,}314$). The checkpoint that produced this figure is \emph{not} part of the released artefact set, and the figure is \emph{not reproducible} from the public training pipeline: \texttt{training/train\_tft\_v3.py} re-run on Vertex AI with the same hyperparameters converges to $51.79\%$. Three explanations are compatible with the gap: (i) train/test leakage in the pre-Vertex setup eliminated by the Vertex container; (ii) a different test-set definition in the development logs; (iii) a different quantile-to-direction decision rule. Without the original checkpoint we cannot discriminate, but we disclose the figure in line with reproducibility-by-checklist practice~\cite{ref:pineau_repro,ref:bouthillier_repro}, retire it as a headline, and adopt the Vertex-reproducible $51.79\%$ as the reference TFT v3 result.}}
\smallskip

\textbf{Per-stock heterogeneity caveat.} Across the ten non-TFT baselines, mean per-ticker accuracy stays in $[49.9\%, 52.9\%]$ and per-ticker median in $[48.1\%, 53.0\%]$, with per-ticker IQR of $7$--$13$ pp depending on the baseline. The maximum single-ticker accuracy reached by any of the ten is $68.5\%$ (always-DOWN on the structurally bearish CARTHAGE CEMENT, $165$ DOWN / $76$ UP over $241$ trading days, a per-stock class imbalance the constant rule trivially exploits); the comparable maxima for logistic regression, random forest, and the standard Transformer sit in the $67$--$68\%$ band on similarly imbalanced tickers. The fraction of tickers below $58.1\%$ (historical always-DOWN majority on the full $2016$--$2025$ corpus) is $\geq 74\%$ for every baseline tested --- $74.4\%$ for the standard Transformer, $79.1\%$ for the simple LSTM, $81.4\%$ for always-DOWN itself, $86.0\%$ for momentum, $88.4\%$ for MA$_{20}$ crossover --- meaning that no baseline meaningfully clears the historical bearish benchmark on more than a quarter of the actively trading tickers. The per-ticker breakdown of TFT v3 is shipped at \texttt{results/baselines/multiagent\_offline\_sweep\_predictions.csv} and aggregated via \texttt{scripts/per\_stock\_breakdown.py}.

\subsection{TFT Internal Ablations}\label{sec:tft_ablations}

Two ablations modify exactly one component of v3, all other hyperparameters held fixed: TFT $-$ VSN replaces the static/encoder/decoder softmax-gated weights by uniform $1/N$ weights, isolating the gating contribution of the Variable Selection Network from its Gated Residual Network contribution; TFT $-$ per-stock norm replaces \texttt{GroupNormalizer(groups=[ticker\_id])} by a global \texttt{TorchNormalizer} fitted across all $68$ tickers (the naive baseline against which per-stock normalisation was originally posed to demonstrate gain). Table~\ref{tab:tft_ablations} summarises the result. Observed deltas vs.\ the reference are $-0.23$~pp ($p = 0.734$) and $-0.39$~pp ($p = 0.672$), both within noise. The original hypothesis that these components carry positive directional skill is neither confirmed nor falsified; the evidence is consistent with both contributing approximately zero accuracy on BVMT 2025 within seed and training-stochasticity noise. Macro-F1 degrades from $0.516$ to $0.503$/$0.471$ as components are removed, consistent with the ablations producing more class-biased classifiers (TFT $-$ per-stock norm predicts UP on $\sim\!60\%$ of rows vs.\ $49.84\%$ slice prior).

\begin{table}[htbp]
\centering
\caption{TFT internal-ablation suite on the harmonized 2025 slice. Each ablation row trains the v3 architecture with exactly one component modified versus the reference; other hyperparameters are held fixed (encoder length 30, hidden size 64, dropout 0.3, pinball quantile loss, 3-epoch warmup). McNemar $p$ is two-sided versus the reference.}
\label{tab:tft_ablations}
\scriptsize
\setlength{\tabcolsep}{4pt}
\begin{tabular}{lcccc}
\toprule
\textbf{Variant} & \textbf{Component modified} & \textbf{Acc.} & \textbf{Macro-F1} & \textbf{McN.\ $p$} \\
\midrule
\textbf{TFT v3 (reference)}   & ---                        & $\mathbf{51.79\%}$ & $\mathbf{0.516}$ & ---     \\
TFT $-$ VSN                   & uniform variable weights   & $51.56\%$ & $0.503$ & $0.734$ \\
TFT $-$ per-stock norm        & global \texttt{TorchNorm.} & $51.40\%$ & $0.471$ & $0.672$ \\
\midrule
Always-UP (slice majority)    & ---                        & $49.84\%$ & $0.333$ & $0.054$ \\
\bottomrule
\end{tabular}
\end{table}

We retain both components in the released architecture for their robustness role (Section~\ref{sec:security}): per-stock target standardisation isolates one stock's price distribution from another's, which the global normaliser does not; the VSN's per-variable importance weights produce an audit trail that the EU AI Act-aligned XAI layer consumes directly. On this dataset they impose no measurable accuracy cost.

\textbf{Quantile calibration.} The pinball quantile loss minimises a proper scoring rule but does not by itself guarantee that the predicted intervals are empirically well-calibrated~\cite{ref:kuleshov,ref:confidence_calibration}. The deployed v3 predicts the $0.10$, $0.50$, and $0.90$ quantiles of the seven-day fractional-return distribution; if the model were calibrated, the realised return would lie inside $[\hat{Q}_{0.1}, \hat{Q}_{0.9}]$ in $80\%$ of test rows. Confidence as defined in Section~\ref{sec:methodology}.C reads from the central interval width, so a miscalibrated width translates directly into a miscalibrated gating decision. We do not validate empirical $80\%$ coverage in this paper because the slice-level coverage diagnostic is a separate validation step from the directional-accuracy comparison that motivates Table~\ref{tab:baselines}; we list it as a limitation in Section~\ref{sec:conclusion}. The per-row CSV ships all three quantiles, enabling the coverage calculation without re-training.

\subsection{Multi-Agent Ablation (Offline-Realisable Subset)}\label{sec:multiagent_ablation}

Of the nine planned configurations, three (A1, A4, A8) are realisable from the offline TFT v3 predictions plus raw price history; six (A2, A3, A5, A6, A7, A9) require the deployed PostgreSQL database (\texttt{news\_articles} pre-scores + \texttt{financial\_ratios}) and the trained XGBoost fundamental checkpoint, none of which is shipped as offline data. For A4 (TFT + Graph) we replace the trained GAT by a structural proxy: $60$-day Pearson correlation across all $68$ tickers, top-$10$ neighbours with $|\rho| \geq 0.30$, graph score $=$ signed-correlation-weighted neighbour TFT v3 predictions in $[-1, +1]$, confidence $= \overline{|\rho|}$ of selected neighbours. Coverage is $69\%$ of rows ($3{,}851/5{,}582$), of which $61.5\%$ ($3{,}435$) clear the deployed orchestrator gate ($\text{conf}\!\ge\!0.20$ and $|\text{score}|\!\ge\!0.08$).

Table~\ref{tab:multiagent_ablation} summarises the result. A1 $=$ A4 $=$ A8 $=$ $51.79\%$ (zero flips, McNemar $p = 1.000$). With deployed weights ($w_\text{tech}\!=\!0.34$, $w_\text{graph}\!=\!0.20$, orchestrator thresholds $0.55/0.45$), the binary technical signal mathematically dominates the continuous graph signal: the combined score lies in $[0.63, 1.0]$ when technical$=$UP and in $[0, 0.37]$ when technical$=$DOWN, so the graph can never flip the decision. A weight-sensitivity sweep (Table~\ref{tab:weight_sensitivity}) confirms the graph begins to flip predictions only at $w_\text{graph} \geq 0.50$ and never significantly differs from A1. The offline graph proxy carries no detectable directional skill above TFT v3 alone, regardless of weighting strategy.

\begin{table}[htbp]
\centering
\caption{Multi-agent ablation on the harmonized 2025 slice. Three configurations are realisable from offline artefacts; six require the deployed PostgreSQL database (sentiment pre-scores, financial ratios) and the trained XGBoost checkpoint. McNemar $p$ is two-sided versus A1.}
\label{tab:multiagent_ablation}
\scriptsize
\setlength{\tabcolsep}{4pt}
\begin{tabular}{llccc}
\toprule
\textbf{Config} & \textbf{Setup} & \textbf{Acc.} & \textbf{$\Delta$ A1} & \textbf{McN.\ $p$} \\
\midrule
A1 & TFT only (Vertex reference)                 & $\mathbf{51.79\%}$ & ---             & ---       \\
A4 & + Graph (proxy, deployed weights)           & $51.79\%$ & $\pm 0.00$ pp   & $1.000$   \\
A8 & A4 $-$ gating (graph always used)           & $51.79\%$ & $\pm 0.00$ pp   & $1.000$   \\
\midrule
A2 & + Fundamental                                & \multicolumn{3}{l}{\textit{DB-dependent}} \\
A3 & + Sentiment                                  & \multicolumn{3}{l}{\textit{DB-dependent}} \\
A5 & + Fundamental + Sentiment                    & \multicolumn{3}{l}{\textit{DB-dependent}} \\
A6 & Seven-agent (deployed)                       & \multicolumn{3}{l}{\textit{DB-dependent}} \\
A7 & A6 $-$ data-quality rules                    & \multicolumn{3}{l}{\textit{DB-dependent}} \\
A9 & A6 $-$ weight redistribution                 & \multicolumn{3}{l}{\textit{DB-dependent}} \\
\midrule
\multicolumn{2}{l}{Graph-only (no TFT)}            & $50.75\%$ & $-1.04$ pp      & $0.271$   \\
\bottomrule
\end{tabular}
\end{table}

\begin{table}[htbp]
\centering
\caption{Weight-sensitivity sweep for A4: varying the orchestrator's graph weight $w_\text{graph}$ with $w_\text{tech} = 1 - w_\text{graph}$, gate active. ``Flips'' counts per-row predictions where A4 disagrees with A1.}
\label{tab:weight_sensitivity}
\scriptsize
\setlength{\tabcolsep}{4pt}
\begin{tabular}{ccccc}
\toprule
$w_\text{graph}$ & \textbf{Acc.} & \textbf{95\% CI} & \textbf{McN.\ vs A1} & \textbf{Flips} \\
\midrule
$0.00$ (TFT only)     & $51.79\%$ & $[50.48, 53.12]$ & ---     & $0$       \\
$0.20$ (deployed)     & $51.79\%$ & $[50.48, 53.12]$ & $1.000$ & $0$       \\
$0.37$ (re-norm.)      & $51.79\%$ & $[50.48, 53.12]$ & $1.000$ & $0$       \\
$0.50$                 & $51.40\%$ & $[50.09, 52.74]$ & $0.476$ & $866$     \\
$0.65$                 & $50.86\%$ & $[49.59, 52.17]$ & $0.176$ & $1{,}422$ \\
\bottomrule
\end{tabular}
\end{table}

\textbf{Falsifiable hypotheses, status on the offline subset.} (i) If $|\,\text{A6} - \text{A1}\,|\!<\!1$ pp the multi-agent layer is performance-neutral; for the offline-realisable subset \{A1, A4, A8\} the gap is $0.00$ pp, consistent with the multi-agent layer contributing via robustness and auditability rather than accuracy. (ii) A8 $=$ A4 exactly (no flips), so confidence gating contributes nothing to accuracy in this 2-agent ensemble; the gating-defense argument of Section~\ref{sec:security} survives as an information-theoretic property (cannot be worse than A4) but is empirically inactive on this dataset. (iii) Per-segment breakdown (banks $\mid$ insurance $\mid$ non-banks) remains DB-dependent.

\subsection{Economic Backtest of Baselines}\label{sec:economic}

Directional accuracy above the random floor does not guarantee exploitability: spread, volume, and execution friction can wipe the signal out, and on a thin market with a structural directional drift even a high accuracy can lose to buy-and-hold. The backtest-overfitting literature in quantitative finance has long documented how multiple-testing inflation and out-of-sample over-fitting produce strategies that look statistically significant in development and disappear in deployment~\cite{ref:bailey_pseudo,ref:harvey_presidential}. We reconstruct a pseudo-PnL series for every baseline by converting each row's prediction into a $7$-day position opened at \texttt{close}$_t$ and closed at \texttt{close}$_{t+7}$, equal-weighted across active tickers, charging a $50$ bps round-trip cost (conservative end of BVMT bid-ask + brokerage). We report the deflated Sharpe ratio of~\cite{ref:deflated_sharpe} with $k\!=\!9$ trials. Table~\ref{tab:economic} reports the long-only NET regime (always-DOWN is flat under long-only and is omitted).

\begin{table}[htbp]
\centering
\caption{Long-only economic backtest of 8 baselines on the 2025 test slice, NET of $50$ bps round-trip cost, $r_f = 8\%$ (BCT TMM proxy). DSR$_p$ is the deflated-Sharpe $p$-value with $k = 9$; bold $p$-values clear $0.05$. Sorted by Sharpe.}
\label{tab:economic}
\scriptsize
\setlength{\tabcolsep}{4pt}
\begin{tabular}{lcccc}
\toprule
\textbf{Baseline} & \textbf{Cum.} & \textbf{Sharpe} & \textbf{Max DD} & \textbf{DSR$_p$} \\
\midrule
Standard Transformer            & $+38.7\%$ & $\mathbf{4.40}$ & $\mathbf{-4.6\%}$ & $\mathbf{0.002}$ \\
Simple LSTM                     & $+38.0\%$ & $3.80$ & $-12.1\%$ & $\mathbf{0.014}$ \\
Logistic regression (L2)        & $\mathbf{+45.8\%}$ & $3.69$ & $-12.6\%$ & $\mathbf{0.021}$ \\
Momentum 5/20                   & $+47.8\%$ & $3.40$ & $-12.7\%$ & $\mathbf{0.038}$ \\
MA$_{20}$ crossover             & $+38.1\%$ & $2.54$ & $-15.5\%$ & $0.152$ \\
XGBoost                          & $+20.4\%$ & $1.66$ & $-17.0\%$ & $0.418$ \\
Always-UP                        & $+28.0\%$ & $1.24$ & $-27.3\%$ & $0.576$ \\
Random forest                    & $+4.2\%$  & $-0.75$ & $-10.1\%$ & $0.982$ \\
\bottomrule
\end{tabular}
\end{table}

Three observations: (i) the 2025 test year had a strong positive drift (always-UP earns $+28\%$ cumulative net) so the binding benchmark is always-UP, not always-DOWN; (ii) four baselines clear the deflated-Sharpe bar (Standard Transformer, LSTM, LogReg L2, Momentum 5/20) but do so on a market with $+28\%$ realised drift, so the test year is favourable to long-biased strategies; (iii) the deep baselines concentrate their modest directional skill on days where the realised return is large, producing a markedly better Sharpe ($4.40$) and drawdown ($-4.6\%$) than logistic regression despite indistinguishable pooled accuracy --- a behaviour that raw accuracy alone misses. TFT v3's economic profile on the harmonized slice is shipped at \texttt{results/baselines/vertex\_slice\_harmonized\_predictions.csv} and falls in the LSTM/Transformer range, consistent with the directional-accuracy ranking.

\subsection{Five Observations}

\begin{itemize}
\item \textbf{Trivial baselines are statistically indistinguishable.} Pairwise McNemar $p$-values among always-UP / always-DOWN / MA$_{20}$ crossover are all above $0.70$. The four trivial / trend rows form a flat $\sim\!50\%$ floor.
\item \textbf{Tabular ML adds modest but real signal.} Logistic regression beats the trivial floor at $52.79\%$ with the highest macro-F1 ($0.524$); McNemar against random forest gives $p = 1.6 \times 10^{-3}$.
\item \textbf{Sequential deep baselines do not break the $\sim\!53\%$ ceiling.} A vanilla two-layer LSTM and a standard 2-block Transformer reach $52.69\%$ and $52.45\%$ respectively on the same 60-day $\times$ 15-feature input; CIs overlap with logistic regression.
\item \textbf{TFT v3 sits inside the same $50$--$53\%$ band.} The reference TFT v3 ($51.79\%$) is statistically indistinguishable from the strongest non-TFT baselines and from always-UP on the same slice. Two TFT ablations move the variant by $\leq 0.4$~pp.
\item \textbf{Thin-market implication.} No model in the panel produces a directional accuracy meaningfully separating from the slice's $50\%$ floor. The non-reproducible $77.4\%$ figure would have implied a skill exceeding the typical $55$--$65\%$ TFT ceiling on developed-market data~\cite{IEEEhowto:kopka}; the reproducible $51.79\%$ is consistent with the literature prior on emerging-market forecasting~\cite{ref:bekaert_harvey,ref:harvey_multiple,ref:harvey_presidential}.
\end{itemize}

\section{Robustness, Auditability, and EU AI Act Alignment}\label{sec:security}

The analyses in this section are \emph{architectural}: we describe threat models and audit-trail properties that follow from the design choices, without running end-to-end poisoning experiments~\cite{ref:data_poisoning}. An empirical adversarial study including a SHAP-fooling baseline~\cite{ref:slack_fooling} is identified as future work.

\textbf{Thin-market data poisoning surface.} An adversary with modest capital can manipulate a single illiquid stock to affect global statistics. Per-stock target normalization isolates each stock $i$'s $(\mu_i, \sigma_i)$ from manipulation of stock $j$; per-encoder-window feature scaling additionally resets local statistics every $L=30$ days. The attack surface is reduced to direct signal injection within a single window of a single stock, and confidence gating ($w_\text{gated} = w_\text{orig} \cdot \text{conf}/\tau$) provides a secondary defense by down-weighting low-confidence outputs in the ensemble. The ablation evidence (Section~\ref{sec:directional_acc}) shows these defenses impose no measurable accuracy cost on BVMT 2025.

\textbf{XAI as audit trail.} Every prediction produces a human-readable explanation describing per-agent contributions, feature attributions (SHAP for XGBoost, attention weights for TFT), and confidence breakdowns. These audit trails \emph{provide the technical building blocks} for the transparency obligations of EU AI Act Articles 13--14~\cite{ref:eu_ai_act}, while explicitly acknowledging that full conformity also depends on governance, risk management (Art. 9), data governance (Art. 10), technical documentation (Annex IV), post-market monitoring (Art. 72), and where applicable conformity assessment by a notified body (Art. 43). Empirical assumptions remain unvalidated: SHAP attributions are not formally tested for faithfulness on BVMT data, and adversarial robustness of explanations~\cite{ref:slack_fooling,ref:alvarez_melis,ref:kumar_shap} is left to future work. We therefore phrase every regulatory claim as \emph{aligned with} or \emph{designed to support} the EU AI Act, not as ``satisfying'' or ``conforming to'' it.

\section{Reproducibility and Artifact Availability}\label{sec:repro}

The released artefacts support a one-command reproduction of every numerical claim in this paper. (i) \textbf{Training pipeline}: the Vertex-AI container (\texttt{vertex/Dockerfile}, \texttt{vertex/cloudbuild.yaml}) and the ablation runner (\texttt{training/tft\_ablations/runner.py}) re-train the TFT v3 reference and the two internal ablations (TFT $-$ VSN, TFT $-$ per-stock norm) bit-for-bit from \texttt{BVMT\_SEED}, in $71$--$161$ minutes on a single NVIDIA T4. (ii) \textbf{Per-row predictions}: \texttt{results/tft\_ablations/vertex/baseline\_test\_predictions.csv} contains $(Q_{0.1}, Q_{0.5}, Q_{0.9}, r_7^\text{true})$ for every test row of the v3 reference and its two ablations. (iii) \textbf{Harmonized panel}: \texttt{results/baselines/vertex\_slice\_harmonized\_predictions.csv} contains per-row predictions for all ten non-TFT baselines plus the three TFT variants on the canonical $N=5{,}582$ slice; the matching \texttt{*.json} carries every number in Table~\ref{tab:baselines}. (iv) \textbf{Multi-agent sweep}: \texttt{scripts/multiagent\_offline\_sweep.py} reproduces Tables~\ref{tab:multiagent_ablation} and~\ref{tab:weight_sensitivity} in $\sim\!60$~s on commodity CPU. (v) \textbf{Sentiment validation tooling}: \texttt{scripts/sample\_news\_for\_labelling.py} and \texttt{scripts/evaluate\_sentiment\_labels.py} produce the 200-article expert-labelling sample and the Cohen's $\kappa$ evaluation report listed in the limitations.

The dataset is sourced from PostgreSQL tables populated from public BVMT price feeds (\texttt{histo\_cotation\_*.csv}) and from the financial-statement and news scrapers in \texttt{scripts/scrape\_ilboursa.py}; the post-cleaning TFT feature file (\texttt{data/features/tft\_features.csv}, $105{,}191$ rows, $13$~MB) is released. Raw news articles and per-ticker financial-ratio rows depend on a database snapshot that is not yet published; this is the dominant remaining reproducibility limitation (Section~\ref{sec:conclusion}).

\section{Conclusion}\label{sec:conclusion}

We present a reproducibility-first multi-component decision pipeline for thin-market equity forecasting on BVMT 2016--2025. The Vertex-reproducible TFT v3 reference reaches $51.79\%$ directional accuracy (95\% CI $[50.43, 53.08]$) on the harmonized 2025 test slice, statistically indistinguishable from ten classical and deep baselines and from always-UP. Two internal ablations move accuracy by $\leq 0.4$~pp ($p > 0.65$), providing no significant evidence for or against the original architectural hypotheses. A multi-agent offline-realizable subset (A1, A4, A8) shows the graph signal cannot flip the binary technical decision under the deployed orchestrator weights; six remaining configurations are DB-dependent. An earlier internal estimate of $77.4\%$ is documented as non-reproducible and retained as a transparency disclosure.

\subsection{Limitations}

\begin{itemize}
\item \textbf{Reproducibility of the historical $77.4\%$ figure}: the dominant limitation. Without the original checkpoint we cannot fully discriminate between train/test leakage, a different test-set definition, and a different decision rule.
\item \textbf{Multi-agent ablation coverage}: only 3 of 9 configurations realisable offline; 6 require the deployed PostgreSQL database (news pre-scores, financial ratios) and the trained fundamental checkpoint.
\item \textbf{Sentiment validation}: the deployed sentiment model has not been validated on the BVMT corpus; per-article Cohen's $\kappa$ on a 200-article expert sample is shipped tooling but not yet executed.
\item \textbf{Arabic coverage}: $\sim\!15\%$ of news articles are Arabic and scored by a French-only model; their sentiment values should be treated as missing.
\item \textbf{Quantile calibration}: $80\%$ prediction intervals are not yet validated against empirical coverage rates.
\item \textbf{Coverage of the 2025 test set}: the harmonized slice has $5{,}582$ row-level predictions over $43$ of the $68$ tickers; performance on the absent $25$ is unknown.
\end{itemize}

\subsection{Future Work}

(i) Replay A2, A3, A5, A6, A7, A9 once the database snapshot is published. (ii) Empirical poisoning study validating the architectural defenses of Section~\ref{sec:security}. (iii) Quantile-calibration validation against empirical $80\%$ coverage rates. (iv) AraBERT-based Arabic sentiment integration to lift coverage from $85\%$ to $100\%$. (v) Per-stock breakdown table with bootstrap CI and lift over always-UP on the harmonized slice (CSV shipped, table omitted for space).

\begin{thebibliography}{99}

\bibitem{IEEEhowto:kopka}
B. Lim, S. Ö. Arik, N. Loeff, and T. Pfister, ``Temporal fusion transformers for interpretable multi-horizon time series forecasting,'' \textit{International Journal of Forecasting}, vol. 37, no. 4, pp. 1748--1764, 2021.

\bibitem{ref:vaswani}
A. Vaswani et al., ``Attention is all you need,'' in \textit{Advances in Neural Information Processing Systems (NeurIPS)}, 2017.

\bibitem{ref:lstm}
S. Hochreiter and J. Schmidhuber, ``Long short-term memory,'' \textit{Neural Computation}, vol. 9, no. 8, pp. 1735--1780, 1997.

\bibitem{ref:gat}
P. Veličković et al., ``Graph attention networks,'' in \textit{ICLR}, 2018.

\bibitem{ref:shap}
S. M. Lundberg and S.-I. Lee, ``A unified approach to interpreting model predictions,'' in \textit{NeurIPS}, 2017.

\bibitem{ref:eu_ai_act}
European Parliament and Council, ``Regulation (EU) 2024/1689 on Artificial Intelligence,'' \textit{Official Journal of the European Union}, 2024.

\bibitem{ref:explainability}
Z. C. Lipton, ``The mythos of model interpretability,'' \textit{Queue}, vol. 16, no. 3, pp. 31--57, 2018.

\bibitem{ref:financial_nlp}
M. Rguibi, S. Ben Amar, and M. H. Kacem, ``French financial sentiment analysis: A case study of domain-specific language models,'' in \textit{Proc. Conference on Computational Linguistics and Intelligent Text Processing}, 2023.

\bibitem{ref:multi_agent}
T. Lux and M. Marchesi, ``Scaling and criticality in a stochastic multi-agent model of a financial market,'' \textit{Nature}, vol. 397, pp. 498--500, 1999.

\bibitem{ref:data_poisoning}
B. Nelson et al., ``Misleading learners: Co-opting your spam filter,'' in \textit{Machine Learning in Cyber-Trust}, pp. 17--51, 2013.

\bibitem{ref:ensemble}
L. I. Kuncheva, \textit{Combining Pattern Classifiers: Methods and Algorithms}, Wiley-Interscience, 2004.

\bibitem{ref:deflated_sharpe}
D. H. Bailey and M. López de Prado, ``The deflated Sharpe ratio,'' \textit{Journal of Portfolio Management}, vol. 40, no. 5, pp. 94--107, 2014.

\bibitem{ref:informer}
H. Zhou et al., ``Informer: Beyond efficient transformer for long sequence time-series forecasting,'' in \textit{AAAI}, 2021.

\bibitem{ref:autoformer}
H. Wu et al., ``Autoformer: Decomposition transformers with auto-correlation for long-term series forecasting,'' in \textit{NeurIPS}, 2021.

\bibitem{ref:fedformer}
T. Zhou et al., ``FEDformer: Frequency enhanced decomposed transformer,'' in \textit{ICML}, 2022.

\bibitem{ref:dlinear}
A. Zeng et al., ``Are transformers effective for time series forecasting?'' in \textit{AAAI}, 2023.

\bibitem{ref:patchtst}
Y. Nie et al., ``A time series is worth 64 words: Long-term forecasting with transformers,'' in \textit{ICLR}, 2023.

\bibitem{ref:timesnet}
H. Wu et al., ``TimesNet: Temporal 2D-variation modeling for general time series analysis,'' in \textit{ICLR}, 2023.

\bibitem{ref:nhits}
C. Challu et al., ``N-HiTS: Neural hierarchical interpolation for time series forecasting,'' in \textit{AAAI}, 2023.

\bibitem{ref:itransformer}
Y. Liu et al., ``iTransformer: Inverted transformers are effective for time series forecasting,'' in \textit{ICLR}, 2024.

\bibitem{ref:timemixer}
S. Wang et al., ``TimeMixer: Decomposable multiscale mixing for time series forecasting,'' in \textit{ICLR}, 2024.

\bibitem{ref:fingat}
Y.-L. Hsu, Y.-C. Tsai, and C.-T. Li, ``FinGAT: Financial graph attention networks,'' \textit{IEEE TKDE}, 2021.

\bibitem{ref:mansf}
R. Sawhney et al., ``Deep attentive learning for stock movement prediction from social media text and company correlations,'' in \textit{EMNLP}, 2020.

\bibitem{ref:bekaert_harvey}
G. Bekaert and C. R. Harvey, ``Time-varying world market integration,'' \textit{Journal of Finance}, vol. 50, no. 2, pp. 403--444, 1995.

\bibitem{ref:slack_fooling}
D. Slack et al., ``Fooling LIME and SHAP: Adversarial attacks on post hoc explanation methods,'' in \textit{AIES}, 2020.

\bibitem{ref:alvarez_melis}
D. Alvarez-Melis and T. S. Jaakkola, ``On the robustness of interpretability methods,'' \textit{ICML Workshop on Human Interpretability in ML}, 2018.

\bibitem{ref:kumar_shap}
I. E. Kumar et al., ``Problems with Shapley-value-based explanations as feature importance measures,'' in \textit{ICML}, 2020.

\bibitem{ref:harvey_multiple}
C. R. Harvey, Y. Liu, and H. Zhu, ``\ldots and the cross-section of expected returns,'' \textit{Review of Financial Studies}, vol. 29, no. 1, pp. 5--68, 2016.

\bibitem{ref:veale_aiact}
M. Veale and F. Zuiderveen Borgesius, ``Demystifying the Draft EU Artificial Intelligence Act,'' \textit{Computer Law Review International}, vol. 22, no. 4, pp. 97--112, 2021.

\bibitem{ref:mokander_audit}
J. Mökander et al., ``Auditing large language models: A three-layered approach,'' \textit{AI and Ethics}, vol. 4, pp. 1085--1115, 2023.

\bibitem{ref:arabert}
W. Antoun, F. Baly, and H. Hajj, ``AraBERT: Transformer-based model for Arabic language understanding,'' in \textit{OSACT4 at LREC}, 2020.

\bibitem{ref:yang_fingpt}
H. Yang, X.-Y. Liu, and C. D. Wang, ``FinGPT: Open-source financial large language models,'' \textit{NeurIPS Workshop on Instruction Tuning}, 2023.

\bibitem{ref:lopez_lira_chatgpt}
A. Lopez-Lira and Y. Tang, ``Can ChatGPT forecast stock price movements? Return predictability and large language models,'' \textit{SSRN}, 2023.

\bibitem{ref:pineau_repro}
J. Pineau et al., ``Improving reproducibility in machine learning research,'' \textit{JMLR}, vol. 22, no. 164, pp. 1--20, 2021.

\bibitem{ref:bouthillier_repro}
X. Bouthillier, C. Laurent, and P. Vincent, ``Unreproducible research is reproducible,'' in \textit{ICML}, 2019.

\bibitem{ref:lipton_troubling}
Z. C. Lipton and J. Steinhardt, ``Troubling trends in machine learning scholarship,'' \textit{Queue}, vol. 17, no. 1, pp. 1--33, 2019.

\bibitem{ref:bailey_pseudo}
D. H. Bailey, J. M. Borwein, M. López de Prado, and Q. J. Zhu, ``Pseudo-mathematics and financial charlatanism,'' \textit{Notices of the AMS}, vol. 61, no. 5, pp. 458--471, 2014.

\bibitem{ref:harvey_presidential}
C. R. Harvey, ``Presidential address: The scientific outlook in financial economics,'' \textit{Journal of Finance}, vol. 72, no. 4, pp. 1399--1440, 2017.

\bibitem{ref:bergmeir_cv_ts}
C. Bergmeir and J. M. Benítez, ``On the use of cross-validation for time series predictor evaluation,'' \textit{Information Sciences}, vol. 191, pp. 192--213, 2012.

\bibitem{ref:kuleshov}
V. Kuleshov, N. Fenner, and S. Ermon, ``Accurate uncertainties for deep learning using calibrated regression,'' in \textit{ICML}, 2018.

\bibitem{ref:confidence_calibration}
C. Guo, G. Pleiss, Y. Sun, and K. Q. Weinberger, ``On calibration of modern neural networks,'' in \textit{ICML}, 2017.

\end{thebibliography}

\balance

\end{document}
